\chapter{Analisis}
\label{chap:metodologi}

\section{Analisis Masalah}
\label{sec:gambaran-umum}

Tugas akhir ini bertujuan untuk merancang, mengimplementasikan, dan membandingkan dua pendekatan algoritmik dalam menyelesaikan permasalahan \textit{Colored Queens}: pendekatan berbasis pencarian sistematis menggunakan \textit{Backtracking} yang diperkuat dengan \textit{Forward Checking} dan AC-3 sebagaimana dibahas pada Bagian \ref{subsubsec:backtracking}, serta pendekatan berbasis \textit{metaheuristic} menggunakan \textit{Particle Swarm Optimization} (PSO) diskrit sebagaimana dibahas pada Bagian \ref{subsubsec:pso}. Selain kedua \textit{solver} tersebut, Tugas akhir ini juga mencakup perancangan dan implementasi sebuah aplikasi permainan berbasis antarmuka grafis yang berfungsi sebagai sarana bermain \textit{Colored Queens} dan integrasi \textit{solver}.

Pendekatan pertama adalah \textit{Backtracking} yang diperkuat dengan \textit{Forward Checking} dan AC-3 sebagaimana dibahas pada Bagian \ref{subsubsec:backtracking}. Pendekatan ini bersifat sistematis dan lengkap (\textit{complete}), sehingga selalu menemukan solusi apabila solusi tersebut ada. Pendekatan kedua adalah \textit{Particle Swarm Optimization} (PSO) diskrit sebagaimana dibahas pada Bagian \ref{subsubsec:pso}, yang merupakan \textit{metaheuristic} berbasis populasi yang tidak menjamin kelengkapan solusi namun berpotensi lebih efisien pada ruang pencarian yang besar. Pemilihan kedua pendekatan ini didasarkan pada kontras paradigmanya, sehingga hasil perbandingan dapat memberikan gambaran yang jelas mengenai \textit{trade-off} antara pencarian deterministik dan stokastik pada permasalahan \textit{Colored Queens}.

Sistem yang dibangun terdiri atas tiga komponen utama. Pertama, sebuah \textit{loader} yang bertugas membaca dan mengurai berkas puzzle dalam format JSON ke dalam struktur papan yang dapat diproses oleh \textit{solver}. Kedua, dua buah \textit{solver} yang mengimplementasikan kedua algoritma pencarian secara terpisah dan independen, sehingga keduanya dapat dijalankan dengan masukan yang sama untuk keperluan eksperimen. Ketiga, sebuah antarmuka pengguna yang dibangun menggunakan Java Swing dan JFrame, yang menampilkan papan permainan secara visual serta memungkinkan pengguna untuk berinteraksi dan memainkan \textit{Colored Queens} secara langsung dan mencoba fitur \textit{hint} dan \textit{auto-solve}.

Mengingat aplikasi permainan memerlukan jaminan bahwa solusi selalu tersedia ketika pengguna menggunakan fitur \textit{hint} atau \textit{auto-solve}, \textit{solver} yang diintegrasikan ke dalam aplikasi adalah \textit{backtracking} dengan \textit{forward checking} dan AC-3. Sifat \textit{complete} dari pendekatan ini menjamin ditemukannya solusi apabila solusi tersebut ada, yang merupakan syarat fungsional yang tidak dapat dipenuhi oleh PSO. \textit{Solver} tersebut dijalankan secara asinkron di latar belakang saat pengguna membuka sebuah puzzle, sehingga solusi tersedia tanpa perlu menyimpan seluruh solusi puzzle secara \textit{hardcoded}. \textit{Solver} PSO diskrit tidak diintegrasikan ke dalam aplikasi, melainkan digunakan sebagai pembanding dalam kerangka eksperimen untuk dievaluasi dari segi waktu eksekusi dan tingkat keberhasilan. Alur kerja penelitian secara keseluruhan diilustrasikan pada Gambar \ref{fig:alur-sistem}.

\begin{figure}[!htbp]
	\centering
	\includegraphics[width=0.85\linewidth]{flowchart.png}
	\caption{Alur kerja penelitian secara keseluruhan.}
	\label{fig:alur-sistem}
\end{figure}

Berikut adalah penjelasan setiap proses pada Gambar \ref{fig:alur-sistem}:

\begin{enumerate}
	\item \textbf{Memilih papan yang akan menjadi tes kasus.} Tahap pertama adalah menentukan papan puzzle yang akan digunakan sebagai masukan pengujian. Dataset mencakup papan berukuran $7 \times 7$ hingga $12 \times 12$ dengan masing-masing 250 level, serta papan berukuran $20 \times 20$ dan $30 \times 30$ dengan masing-masing 1 level. Rincian dataset disajikan pada Bagian \ref{subsec:dataset-puzzle}.
	
	\item \textbf{Mengambil data papan dari berkas hasil \textit{scraping}.} Data puzzle diperoleh dari situs \textit{playqueensgame.com} melalui proses \textit{web scraping} otomatis. Untuk setiap level, posisi dan warna setiap sel pada papan diekstrak dari halaman puzzle, kemudian disimpan dalam berkas berformat JSON. Berkas tersebut menjadi sumber data untuk seluruh tahap berikutnya.
	
	\item \textbf{Memodelkan papan permainan sebagai \textit{Constraint Satisfaction Problem} (CSP).} Data papan yang telah dibaca dimodelkan sebagai CSP dengan mendefinisikan variabel (sektor warna), domain (sel-sel dalam setiap sektor), dan kendala (baris, kolom, \textit{adjacency}). Pemodelan CSP secara lengkap dibahas pada Bagian \ref{sec:pemodelan-csp}.
	
	\item \textbf{Memilih algoritma yang akan dijalankan untuk mendapatkan solusi.} Pada tahap ini, salah satu dari dua algoritma dipilih untuk menyelesaikan puzzle. Dalam konteks eksperimen, kedua algoritma dijalankan pada puzzle yang sama untuk keperluan perbandingan.
	
	\item \textbf{Menjalankan algoritma pencarian \textit{backtracking} sampai menemukan solusi.} Apabila algoritma \textit{backtracking} dengan \textit{forward checking} dan AC-3 dipilih, \textit{solver} melakukan pencarian sistematis dengan pemangkasan domain hingga solusi valid ditemukan. Sifat deterministik dan kelengkapan (\textit{completeness}) algoritma ini menjamin ditemukannya solusi apabila solusi tersebut ada. Perancangan algoritma dibahas pada Bagian \ref{sec:analisis-backtracking}.
	
	\item \textbf{Menjalankan algoritma \textit{metaheuristic} PSO hingga menemukan solusi atau sampai suatu \textit{stopping condition}.} Apabila algoritma PSO diskrit dipilih, \textit{solver} mengeksplorasi ruang solusi melalui sejumlah partikel yang berinteraksi dalam \textit{neighborhood}. Pencarian berhenti apabila solusi valid ditemukan, batas iterasi tercapai, atau stagnasi terdeteksi. Perancangan algoritma dibahas pada Bagian \ref{sec:analisis-pso}.
	
	\item \textbf{Solusi valid akan ditampilkan.} Solusi berupa pasangan koordinat penempatan menteri untuk setiap sektor warna dihasilkan oleh solver. Pada konteks aplikasi, solusi dirender secara visual pada papan permainan. Pada konteks eksperimen, solusi beserta metrik kinerja dicatat ke dalam berkas hasil untuk dianalisis lebih lanjut.
	
	\item \textbf{Melakukan perbandingan dan evaluasi.} Metrik dari kedua \textit{solver} diagregasi dan dibandingkan berdasarkan waktu eksekusi, tingkat keberhasilan, dan karakteristik pencarian pada berbagai ukuran papan dan konfigurasi parameter. Skenario pengujian dibahas pada Bagian \ref{sec:skenario-pengujian} dan hasil analisis disajikan pada Bab \ref{chap:hasil}.
\end{enumerate}

\section{Pemodelan Permasalahan \textit{Colored Queens} sebagai CSP}
\label{sec:pemodelan-csp}

Sebagaimana telah dibahas pada Bagian~\ref{sec:csp}, sebuah \textit{Constraint Satisfaction Problem} didefinisikan oleh tiga komponen utama: himpunan variabel, domain masing-masing variabel, dan himpunan kendala yang harus dipenuhi secara bersamaan. Bagian ini menjabarkan secara konkret bagaimana ketiga komponen tersebut didefinisikan dalam konteks permasalahan \textit{Colored Queens}. Pemodelan formal ini menjadi landasan bagi kedua pendekatan penyelesaian yang dirancang pada Bagian~\ref{sec:analisis-backtracking} dan Bagian~\ref{sec:analisis-pso}.

\subsection{Representasi Papan}
\label{subsec:representasi-papan}

Papan permainan \textit{Colored Queens} dimodelkan sebagai sebuah grid berukuran $n \times n$, di mana $n$ menyatakan jumlah baris, jumlah kolom, sekaligus jumlah sektor warna yang ada pada papan. Setiap sel pada papan diidentifikasi oleh pasangan koordinat $(\textit{baris}, \textit{kolom})$ dengan indeks berbasis nol, sehingga sel pojok kiri atas memiliki koordinat $(0, 0)$ dan sel pojok kanan bawah memiliki koordinat $(n-1, n-1)$. Setiap sel memiliki atribut warna yang menentukan sektor mana sel tersebut termasuk.

Sektor warna didefinisikan sebagai himpunan sel-sel yang memiliki warna sama. Pada papan yang valid, terdapat tepat $n$ sektor warna, dan setiap sel pada papan termasuk dalam tepat satu sektor warna. Bentuk dan luas sektor warna dapat berbeda-beda tergantung pada tata letak papan yang diberikan, sebagaimana telah diperlihatkan pada contoh papan di Bagian~\ref{subsec:colored-queens-definisi}.

\subsection{Definisi Variabel}
\label{subsec:definisi-variabel}

Mengacu pada kerangka CSP yang dijelaskan pada Bagian~\ref{sec:csp}, variabel dalam permasalahan \textit{Colored Queens} didefinisikan sebagai sektor warna pada papan. Secara formal, diberikan himpunan warna $C = \{c_1, c_2, \dots, c_n\}$, maka himpunan variabel didefinisikan pada Persamaan~\ref{eq:def-variabel}.

\begin{equation}
	X = \{X_{c_1}, X_{c_2}, \dots, X_{c_n}\}
	\label{eq:def-variabel}
\end{equation}

di mana setiap variabel $X_{c_i}$ merepresentasikan pilihan posisi penempatan menteri untuk sektor warna $c_i$.

Pemilihan sektor warna sebagai variabel, bukan baris atau kolom seperti pada N-Queens klasik, didasarkan pada struktur kendala \textit{Colored Queens} itu sendiri. Aturan permainan mengharuskan tepat satu menteri ditempatkan pada setiap sektor warna (kendala \textit{one queen per color}), sehingga setiap sektor warna secara alami berkorespondensi dengan satu variabel keputusan. Dengan pemodelan ini, himpunan variabel memiliki kardinalitas yang sama dengan ukuran papan $n$, dan setiap penugasan lengkap (\textit{complete assignment}) secara otomatis memenuhi kendala jumlah menteri per warna.

\subsection{Definisi Domain}
\label{subsec:definisi-domain}

Domain dari setiap variabel $X_{c_i}$ adalah himpunan semua sel yang termasuk dalam sektor warna $c_i$, sebagaimana didefinisikan pada Persamaan~\ref{eq:def-domain}.

\begin{equation}
	D_{c_i} = \{(r, k) \mid \text{sel } (r, k) \text{ memiliki warna } c_i\}
	\label{eq:def-domain}
\end{equation}

Dengan demikian, ukuran domain setiap variabel dapat berbeda-beda tergantung pada bentuk dan luas sektor warna yang bersangkutan pada papan. Tidak seperti N-Queens klasik yang menggunakan representasi permutasi dengan domain seragam $\{1, 2, \dots, n\}$, domain pada \textit{Colored Queens} bersifat heterogen dan bergantung pada tata letak papan yang diberikan.

Selama proses pencarian, status kevalidan setiap sel dalam domain bersifat dinamis. Sel dapat dipangkas dari domain (dianggap tidak lagi valid sebagai kandidat penempatan) oleh mekanisme seperti \textit{forward checking} atau propagasi AC-3, dan dapat dipulihkan kembali ke domain ketika algoritma melakukan \textit{backtrack}. Dengan demikian, domain efektif yang dipertimbangkan oleh algoritma pada suatu titik dalam pencarian merupakan himpunan bagian dari domain awalnya.

\subsection{Definisi \textit{Constraint}}
\label{subsec:definisi-constraint}

Sesuai dengan aturan permainan \textit{Colored Queens} yang dijabarkan pada Bagian~\ref{subsec:colored-queens-definisi}, terdapat tiga jenis kendala yang harus dipenuhi secara bersamaan oleh setiap pasangan variabel $(X_{c_i}, X_{c_j})$ dengan $i \neq j$. Ketiga kendala ini bersifat biner dan membentuk \textit{constraint graph} yang lengkap (\textit{complete}), artinya setiap pasang variabel terhubung oleh kendala.

\textbf{Kendala Baris.} Dua menteri tidak boleh ditempatkan pada baris yang sama. Untuk dua sel $(r_1, k_1)$ dan $(r_2, k_2)$ yang ditugaskan kepada variabel $X_{c_i}$ dan $X_{c_j}$, kendala ini dinyatakan pada Persamaan~\ref{eq:kendala-baris}.
\begin{equation}
	r_1 \neq r_2
	\label{eq:kendala-baris}
\end{equation}

\textbf{Kendala Kolom.} Dua menteri tidak boleh ditempatkan pada kolom yang sama, sebagaimana dinyatakan pada Persamaan~\ref{eq:kendala-kolom}.
\begin{equation}
	k_1 \neq k_2
	\label{eq:kendala-kolom}
\end{equation}

\textbf{Kendala \textit{Adjacency}.} Berbeda dari permasalahan N-Queens klasik yang melarang konflik diagonal tak terbatas, \textit{Colored Queens} hanya melarang penempatan dua menteri pada sel yang bertetangga secara langsung dalam delapan arah (horizontal, vertikal, dan diagonal). Dua sel $(r_1, k_1)$ dan $(r_2, k_2)$ dinyatakan bertetangga apabila memenuhi Persamaan~\ref{eq:kendala-adjacency}.
\begin{equation}
	|r_1 - r_2| \leq 1 \quad \text{dan} \quad |k_1 - k_2| \leq 1
	\label{eq:kendala-adjacency}
\end{equation}

Perlu dicatat bahwa tidak ada larangan bagi dua menteri untuk berada pada diagonal yang sama, selama keduanya tidak bertetangga secara langsung. Hal ini merupakan perbedaan struktural yang signifikan dari N-Queens klasik dan menjadi karakteristik utama dari \textit{Colored Queens} sebagaimana dibahas pada Bagian~\ref{subsec:colored-queens-definisi}.

Ketiga kendala di atas dapat digabungkan menjadi sebuah definisi konflik tunggal antara dua sel sebagaimana dinyatakan pada Persamaan~\ref{eq:def-konflik}. Dua sel dinyatakan berkonflik apabila memenuhi salah satu dari kondisi pada persamaan tersebut.

\begin{equation}
	\text{konflik}((r_1, k_1),\, (r_2, k_2)) \iff (r_1 = r_2) \;\lor\; (k_1 = k_2) \;\lor\; (|r_1 - r_2| \leq 1 \;\land\; |k_1 - k_2| \leq 1)
	\label{eq:def-konflik}
\end{equation}

Pemeriksaan konflik melalui formulasi ini akan digunakan secara konsisten oleh kedua algoritma penyelesaian yang dirancang pada bagian-bagian berikutnya.

\subsection{Representasi Solusi}
\label{subsec:representasi-solusi}

Sebuah solusi valid dari permasalahan \textit{Colored Queens} yang dimodelkan sebagai CSP adalah sebuah penugasan lengkap sebagaimana dinyatakan pada Persamaan~\ref{eq:def-penugasan}.

\begin{equation}
	A = \{X_{c_1} \mapsto s_1,\; X_{c_2} \mapsto s_2,\; \dots,\; X_{c_n} \mapsto s_n\}
	\label{eq:def-penugasan}
\end{equation}

di mana $s_i \in D_{c_i}$ untuk setiap $i$, dan tidak ada pasangan $(s_i, s_j)$ dengan $i \neq j$ yang melanggar kendala konflik yang telah didefinisikan pada Persamaan~\ref{eq:def-konflik}.

Penugasan lengkap inilah yang menjadi keluaran dari kedua algoritma yang dirancang dalam tugas akhir ini. Pada algoritma \textit{backtracking}, penugasan dibangun secara inkremental melalui penambahan satu variabel pada satu waktu hingga seluruh variabel mendapat nilai. Pada algoritma PSO, setiap partikel pada setiap iterasi merepresentasikan satu kandidat penugasan lengkap, meskipun belum tentu memenuhi seluruh kendala.

\section{Analisis Algoritma \textit{Backtracking} dengan \textit{Forward Checking} dan AC-3}
\label{sec:analisis-backtracking}

Bagian ini menjelaskan analisis algoritmik dari pendekatan pertama yang digunakan untuk menyelesaikan permasalahan \textit{Colored Queens}, yaitu \textit{backtracking} sistematis yang diperkuat dengan dua mekanisme pemangkasan domain: \textit{forward checking} dan AC-3. Landasan teori ketiga teknik tersebut telah dibahas pada Bagian~\ref{subsubsec:backtracking} dan Bagian~\ref{subsubsec:ac3}. Analisis disajikan melalui empat diagram alur yang merepresentasikan setiap komponen algoritma secara terpisah, yaitu alur utama rekursif pencarian, prosedur \textit{forward checking}, fungsi \textit{revise} sebagai inti dari konsistensi busur, dan prosedur propagasi AC-3 yang menjalankan propagasi secara iteratif. Rincian implementasi masing-masing komponen dibahas pada Bab~\ref{chap:perancangan}.

\subsection{Alur Utama Algoritma \textit{Backtracking}}
\label{subsec:alur-backtracking}

Alur utama algoritma \textit{backtracking} mengimplementasikan pencarian rekursif \textit{depth-first} yang memproses variabel warna satu per satu sesuai urutan yang telah ditentukan. Pada setiap pemanggilan rekursif, satu menteri ditempatkan untuk satu sektor warna, diikuti oleh pemangkasan domain dan pemanggilan rekursif berikutnya. Apabila pencarian pada suatu cabang gagal, algoritma membatalkan penugasan dan mencoba alternatif lain pada level yang sama. Keseluruhan alur diilustrasikan pada Gambar~\ref{fig:flow-placequeens}.

\begin{figure}[!htbp]
	\centering
	\includegraphics[width=0.95\linewidth]{flowbt.png}
	\caption{Diagram alur prosedur rekursif pencarian \textit{backtracking}.}
	\label{fig:flow-placequeens}
\end{figure}

\subsubsection{Pengecekan Kondisi Terminasi}

Setiap pemanggilan rekursif diawali dengan dua pemeriksaan kondisi terminasi. Pemeriksaan pertama, ``Semua warna sudah ditugaskan?'', menentukan apakah \textit{base case} pencarian telah tercapai. Kondisi ini terpenuhi ketika seluruh sektor warna telah memiliki menteri yang ditempatkan, sehingga prosedur mengembalikan nilai \textit{true} yang menandakan ditemukannya solusi valid. Pemeriksaan kedua, ``Domain warna saat ini kosong?'', memeriksa apakah variabel warna yang sedang diproses tidak memiliki sel valid yang tersisa. Apabila kondisi ini terpenuhi, pencarian pada cabang ini tidak dapat dilanjutkan dan prosedur mengembalikan \textit{false} kepada pemanggilnya sebagai sinyal untuk melakukan \textit{backtrack}.

\subsubsection{Pemilihan dan Penempatan Menteri}

Apabila kedua kondisi terminasi tidak terpenuhi, algoritma masuk ke dalam perulangan utama yang mengiterasi seluruh sel valid dalam domain variabel saat ini. Dari seluruh sel yang tersedia, algoritma memilih satu sel kandidat berikutnya dan menempatkan menteri pada sel tersebut secara tentatif, artinya penempatan ini dapat dibatalkan apabila pencarian pada cabang ini gagal. Penempatan menteri pada sel tertentu secara langsung membatasi pilihan yang tersedia bagi variabel-variabel yang belum ditugaskan, karena setiap sel yang berkonflik dengan posisi tersebut tidak lagi dapat digunakan.

\subsubsection{Penyimpanan \textit{Checkpoint} dan Pemangkasan Domain}

Sebelum menjalankan pemangkasan, algoritma mencatat kondisi domain saat ini sebagai \textit{checkpoint}. Mekanisme ini memungkinkan seluruh perubahan domain yang dilakukan pada langkah berikutnya untuk dibatalkan secara efisien apabila pencarian gagal, tanpa perlu menyalin ulang seluruh struktur domain. Setelah \textit{checkpoint} disimpan, algoritma menjalankan dua tahap pemangkasan domain secara berurutan: \textit{forward checking} yang akan dijabarkan pada Bagian~\ref{subsec:flow-fc}, diikuti oleh propagasi AC-3 yang akan dijabarkan pada Bagian~\ref{subsec:flow-propagate}. Kedua tahap ini secara bersama-sama mengeliminasi sel-sel yang sudah tidak konsisten dari domain seluruh variabel yang belum ditugaskan, mempersempit ruang pencarian sebelum rekursi dilanjutkan.

\subsubsection{Evaluasi Pemangkasan dan Pemanggilan Rekursif}

Setelah pemangkasan selesai, algoritma memeriksa apakah salah satu variabel yang belum ditugaskan kehilangan seluruh kandidat selnya akibat pemangkasan. Apabila kondisi ini terjadi, cabang pencarian ini dipastikan tidak dapat menghasilkan solusi valid, sehingga algoritma langsung melompat ke proses pemulihan tanpa melanjutkan rekursi. Deteksi dini kegagalan ini merupakan salah satu keunggulan utama kombinasi \textit{forward checking} dan AC-3 dibandingkan \textit{backtracking} murni. Apabila tidak ada domain yang kosong, algoritma memanggil dirinya sendiri secara rekursif untuk memproses variabel warna berikutnya. Apabila pemanggilan rekursif tersebut berhasil menemukan solusi, nilai \textit{true} diteruskan kembali ke pemanggil.

\subsubsection{Pemulihan dan Eksplorasi Sel Berikutnya}

Apabila pemanggilan rekursif gagal atau pemangkasan menyebabkan domain variabel lain menjadi kosong, algoritma melakukan \textit{backtrack} dengan memulihkan seluruh domain ke kondisi tepat sebelum penempatan saat ini, menggunakan informasi yang tersimpan pada \textit{checkpoint}. Penempatan menteri yang gagal dibatalkan, dan sel tersebut dianggap tidak layak untuk iterasi saat ini. Algoritma kemudian memeriksa apakah masih terdapat sel valid lain dalam domain yang belum dicoba. Apabila masih ada, perulangan dilanjutkan untuk mencoba sel berikutnya. Apabila seluruh sel telah dicoba tanpa keberhasilan, prosedur mengembalikan \textit{false} kepada pemanggilnya, yang akan memicu \textit{backtrack} pada level rekursi di atasnya.

\subsection{\textit{Forward Checking}}
\label{subsec:flow-fc}

\textit{Forward checking} merupakan tahap pemangkasan pertama yang dijalankan setiap kali seorang menteri ditempatkan pada papan. Sebagaimana dibahas pada Bagian~\ref{subsubsec:backtracking}, mekanisme ini secara reaktif mengeliminasi nilai-nilai domain yang tidak konsisten dengan penugasan yang baru saja dilakukan, sehingga cabang pencarian yang pasti gagal dapat dihindari sedini mungkin. Dalam konteks \textit{Colored Queens}, \textit{forward checking} memangkas seluruh sel yang berkonflik dengan menteri yang baru ditempatkan dari domain variabel-variabel warna yang belum ditugaskan. Alur prosedur ini diperlihatkan pada Gambar~\ref{fig:flow-fc}.

\begin{figure}[!htbp]
	\centering
	\includegraphics[width=0.95\linewidth]{flowfc.png}
	\caption{Diagram alur prosedur \textit{forward checking}.}
	\label{fig:flow-fc}
\end{figure}

\subsubsection{Iterasi atas Warna yang Belum Ditugaskan}

Prosedur dimulai dengan mengiterasi seluruh variabel warna yang belum ditugaskan, yaitu variabel-variabel yang belum memiliki menteri. Iterasi dilakukan secara terurut sesuai urutan pemrosesan variabel yang telah ditetapkan. Apabila seluruh variabel yang belum ditugaskan telah diperiksa, prosedur berakhir dan kendali dikembalikan ke alur utama \textit{backtracking}.

\subsubsection{Iterasi atas Sel dalam Domain Warna Terpilih}

Untuk setiap variabel warna yang diperiksa, prosedur mengiterasi seluruh sel yang masih valid dalam domainnya. Sel yang sebelumnya telah dipangkas oleh penugasan-penugasan sebelumnya tidak diperiksa ulang. Apabila seluruh sel pada domain variabel saat ini telah diperiksa, alur beralih ke variabel warna berikutnya.

\subsubsection{Pemeriksaan Konflik dan Pemangkasan Sel}

Setiap sel yang diperiksa dievaluasi terhadap posisi menteri yang baru ditempatkan. Pemeriksaan konflik mengacu pada definisi kendala \textit{Colored Queens} yang telah dijabarkan pada Bagian~\ref{subsec:definisi-constraint}, yang mencakup kesamaan baris, kesamaan kolom, dan ketetanggaan dalam delapan arah. Apabila sel tidak berkonflik, sel tersebut tetap berada dalam domain dan alur melanjutkan ke sel berikutnya. Apabila sel berkonflik, sel tersebut dihapus dari domain dan pemangkasan ini dicatat agar dapat dipulihkan apabila algoritma melakukan \textit{backtrack} di kemudian hari.

\subsection{Fungsi \textit{Revise}}
\label{subsec:flow-revise}

Fungsi \textit{revise} merupakan inti dari mekanisme konsistensi busur AC-3. Sebagaimana telah dijabarkan pada Bagian~\ref{subsubsec:ac3}, fungsi ini menerima sepasang variabel sebagai parameter, yaitu variabel sumber dan variabel tujuan, kemudian memeriksa apakah setiap nilai dalam domain variabel tujuan memiliki nilai pendamping yang konsisten pada variabel sumber. Nilai yang tidak memiliki pendamping yang konsisten dianggap tidak \textit{arc-consistent} dan dipangkas dari domain. Dalam konteks \textit{Colored Queens}, fungsi ini juga memperhitungkan posisi menteri yang sudah ditempatkan pada variabel-variabel yang telah ditugaskan, sehingga setiap kandidat sel tidak hanya konsisten secara berpasangan, tetapi juga konsisten dengan keseluruhan penugasan parsial yang sedang dibangun. Alur fungsi ini diperlihatkan pada Gambar~\ref{fig:flow-revise}.

\begin{figure}[!htbp]
	\centering
	\includegraphics[width=0.95\linewidth]{flowrevise.png}
	\caption{Diagram alur fungsi \textit{revise}.}
	\label{fig:flow-revise}
\end{figure}

\subsubsection{Inisialisasi dan Iterasi atas Domain Variabel Tujuan}

Fungsi dimulai dengan menetapkan penanda \textit{revised} bernilai \textit{false}. Penanda ini akan menjadi nilai kembalian fungsi yang menunjukkan apakah domain variabel tujuan mengalami perubahan akibat eksekusi. Selanjutnya, fungsi mengiterasi seluruh sel valid dalam domain variabel tujuan. Apabila seluruh sel telah diperiksa, fungsi mengembalikan nilai \textit{revised} sebagai hasil akhir.

\subsubsection{Pemeriksaan Keberadaan \textit{Support} pada Variabel Sumber}

Untuk setiap sel yang diperiksa pada variabel tujuan, fungsi memeriksa apakah terdapat setidaknya satu sel pada variabel sumber yang tidak berkonflik dengannya. Sebuah sel dinyatakan memiliki \textit{support} apabila terdapat minimal satu pasangan nilai yang valid antara kedua variabel. Konflik antara sepasang sel mengacu pada definisi kendala yang telah dijabarkan pada Bagian~\ref{subsec:definisi-constraint}. Apabila tidak ditemukan sel pendukung pada variabel sumber, sel pada variabel tujuan dipastikan tidak \textit{arc-consistent} dan dipangkas dari domain.

\subsubsection{Pemeriksaan Konflik dengan Menteri yang Sudah Ditempatkan}

Apabila sel memiliki \textit{support} dari variabel sumber, fungsi melanjutkan pemeriksaan terhadap seluruh menteri yang telah ditempatkan pada variabel-variabel sebelumnya. Langkah tambahan ini diperlukan karena AC-3 dalam konteks \textit{Colored Queens} tidak hanya menjaga konsistensi antarvariabel yang belum ditugaskan, tetapi juga memastikan bahwa setiap kandidat sel tetap konsisten dengan keseluruhan penugasan parsial yang telah dibangun. Apabila sel berkonflik dengan salah satu menteri yang telah ditempatkan, sel tersebut dipangkas meskipun memiliki \textit{support} pada variabel sumber.

\subsubsection{Pemangkasan Sel Tanpa \textit{Support}}

Pada saat sel terbukti tidak memiliki \textit{support} yang valid, baik karena tidak ditemukan sel pendamping pada variabel sumber maupun karena berkonflik dengan menteri yang telah ditempatkan, sel tersebut dihapus dari domain variabel tujuan. Pemangkasan ini dicatat agar dapat dipulihkan apabila algoritma melakukan \textit{backtrack}, dan penanda \textit{revised} diatur menjadi \textit{true} untuk menandakan bahwa domain mengalami perubahan. Iterasi kemudian dilanjutkan ke sel berikutnya pada domain variabel tujuan.

\subsection{Propagasi AC-3}
\label{subsec:flow-propagate}

Propagasi AC-3 merupakan tahap pemangkasan kedua yang dijalankan setelah \textit{forward checking} selesai. Sebagaimana dijabarkan pada Bagian~\ref{subsubsec:ac3}, AC-3 menegakkan konsistensi busur antarseluruh pasangan variabel yang belum ditugaskan, mengeksploitasi keterkaitan kendala antarvariabel yang tidak terdeteksi oleh \textit{forward checking}. Propagasi dilakukan secara iteratif menggunakan antrean busur: setiap busur diambil, diperiksa konsistensinya melalui fungsi \textit{revise} sebagaimana dibahas pada Bagian~\ref{subsec:flow-revise}, dan apabila terjadi perubahan domain, busur-busur lain yang terkait ditambahkan ke antrean untuk dipropagasi lebih lanjut. Alur prosedur ini diperlihatkan pada Gambar~\ref{fig:flow-propagate}.

\begin{figure}[!htbp]
	\centering
	\includegraphics[width=0.75\linewidth]{flowac3.png}
	\caption{Diagram alur prosedur propagasi AC-3.}
	\label{fig:flow-propagate}
\end{figure}

\subsubsection{Pengecekan Antrean dan Pengambilan Busur}

Prosedur dimulai dengan memeriksa apakah antrean busur masih berisi busur yang belum diproses. Antrean awal diinisialisasi dengan seluruh pasang busur antara variabel-variabel yang belum ditugaskan. Apabila antrean kosong, propagasi telah selesai dan seluruh busur dipastikan konsisten, sehingga prosedur berakhir. Apabila antrean masih berisi busur, satu busur diambil untuk diproses pada langkah berikutnya.

\subsubsection{Pemanggilan Fungsi \textit{Revise}}

Busur yang diambil dari antrean diproses dengan menjalankan fungsi \textit{revise} sebagaimana dijelaskan pada Bagian~\ref{subsec:flow-revise}. Fungsi ini memangkas sel-sel pada domain variabel tujuan yang tidak memiliki \textit{support} pada variabel sumber atau yang berkonflik dengan menteri yang telah ditempatkan, kemudian mengembalikan status perubahan domain. Status ini menentukan langkah berikutnya pada propagasi.

\subsubsection{Evaluasi Hasil Pemangkasan}

Apabila fungsi \textit{revise} tidak menghasilkan perubahan pada domain variabel tujuan, dampak propagasi terhenti pada busur tersebut, dan alur kembali memeriksa antrean. Apabila domain mengalami perubahan, prosedur selanjutnya memeriksa apakah domain variabel tujuan menjadi kosong. Jika demikian, hal ini menandakan terjadinya \textit{domain wipeout}. Pada kondisi ini, propagasi dihentikan dan sinyal kegagalan dikembalikan kepada alur utama \textit{backtracking} agar dapat melakukan \textit{backtrack}.

\subsubsection{Pengembangan Antrean dengan Busur Terkait}

Apabila domain variabel tujuan mengalami perubahan namun tidak menjadi kosong, dampak pemangkasan tersebut dapat memengaruhi konsistensi busur lain yang melibatkan variabel tujuan sebagai variabel sumber. Oleh karena itu, seluruh busur yang berbentuk $(t, k)$, di mana $t$ adalah variabel tujuan saat ini dan $k$ adalah variabel lain yang belum ditugaskan, ditambahkan kembali ke antrean. Mekanisme ini memastikan dampak pemangkasan terpropagasi secara transitif ke seluruh variabel yang saling terhubung. Setelah penambahan selesai, alur kembali ke pemeriksaan antrean untuk melanjutkan iterasi.

\subsection{Studi Kasus: Penyelesaian Papan 7$\times$7}
\label{subsec:studi-kasus-bt}

Untuk memperjelas cara kerja algoritma \textit{backtracking} yang diperkuat dengan \textit{forward checking} dan AC-3, bagian ini menyajikan studi kasus penyelesaian sebuah papan \textit{Colored Queens} berukuran 7$\times$7 yang terdiri atas tujuh sektor warna, yaitu ungu, biru, hijau, oranye, merah, abu-abu, dan kuning. Studi kasus ini menelusuri proses penyelesaian dengan penekanan pada bagaimana kedua mekanisme pemangkasan mengurangi ruang pencarian dan bagaimana \textit{backtrack} terjadi ketika cabang pencarian gagal.

\subsubsection{Penempatan Awal dan \textit{Forward Checking}}

Algoritma memulai pencarian dengan memproses variabel warna pertama, dalam contoh ini warna biru. Sebuah sel valid dipilih dari domain warna biru sebagai posisi penempatan menteri, sebagaimana diperlihatkan pada Gambar~\ref{fig:bt-step20}.

\begin{figure}[!htbp]
	\centering
	\includegraphics[width=0.27\linewidth]{./bt/20.png}
	\caption{Penempatan menteri pertama pada salah satu sel dari domain warna biru.}
	\label{fig:bt-step20}
\end{figure}

Setelah menteri ditempatkan, \textit{forward checking} dijalankan untuk memangkas sel-sel yang berkonflik langsung dengan menteri tersebut dari domain seluruh variabel warna yang belum ditugaskan. Hasil pemangkasan diperlihatkan pada Gambar~\ref{fig:bt-step21}: seluruh sel yang berada pada baris atau kolom yang sama dengan menteri, serta sel-sel yang bertetangga secara langsung dalam delapan arah, ditandai sebagai tidak valid (tanda X merah).

\begin{figure}[!htbp]
	\centering
	\includegraphics[width=0.27\linewidth]{./bt/21.png}
	\caption{Hasil pemangkasan oleh \textit{forward checking} setelah penempatan menteri pertama.}
	\label{fig:bt-step21}
\end{figure}

\subsubsection{Propagasi AC-3}

Setelah \textit{forward checking} selesai, algoritma melanjutkan dengan propagasi AC-3 untuk menegakkan konsistensi busur antarvariabel yang belum ditugaskan. Konsep dasar dari pemeriksaan satu busur (\textit{arc}) diilustrasikan pada Gambar~\ref{fig:bt-step22}: untuk setiap pasang variabel, AC-3 memastikan bahwa setiap nilai dalam domain salah satu variabel memiliki minimal satu nilai pendamping yang konsisten (\textit{support}) pada variabel lainnya.

\begin{figure}[!htbp]
	\centering
	\includegraphics[width=0.75\linewidth]{./bt/22.png}
	\caption{Konsep dasar busur (\textit{arc}) dalam AC-3.}
	\label{fig:bt-step22}
\end{figure}

Dalam konteks \textit{Colored Queens}, setiap variabel warna saling terhubung dengan seluruh variabel warna lainnya, membentuk graf kendala yang lengkap. Antrean awal AC-3 diinisialisasi dengan seluruh pasang busur antara variabel-variabel yang belum ditugaskan, sebagaimana diperlihatkan pada Gambar~\ref{fig:bt-step23}.

\begin{figure}[!htbp]
	\centering
	\includegraphics[width=0.7\linewidth]{./bt/23.png}
	\caption{Graf kendala lengkap dan antrean awal busur AC-3.}
	\label{fig:bt-step23}
\end{figure}

Sebagai ilustrasi pemrosesan satu busur, perhatikan busur Hijau$\rightarrow$Kuning yang menjadi busur pertama dalam antrean. Pada titik ini, domain warna hijau yang masih valid setelah \textit{forward checking} terdiri atas sel-sel $\{(1,2),\,(1,3),\,(1,4),\,(4,4),\,(5,2),\,(5,3),\,(5,4)\}$, sementara domain warna kuning terdiri atas $\{(6,2),\,(6,3)\}$. Algoritma memeriksa setiap sel valid pada domain warna hijau dan menentukan apakah terdapat setidaknya satu sel pada domain warna kuning yang tidak berkonflik dengannya.

Gambar~\ref{fig:bt-step24} menunjukkan pemrosesan sel Hijau(1,2), yang ditandai dengan kotak putih pada papan. Sel ini berkonflik dengan Kuning(6,2) karena berada pada kolom yang sama, tetapi tidak berkonflik dengan Kuning(6,3) karena berada pada baris, kolom, dan ketetanggaan yang berbeda. Dengan demikian, Hijau(1,2) memiliki minimal satu \textit{support} pada domain warna kuning, sehingga sel ini tetap dipertahankan dalam domain.

\begin{figure}[!htbp]
	\centering
	\includegraphics[width=0.27\linewidth]{./bt/24.png}
	\caption{Pemrosesan busur Hijau$\rightarrow$Kuning: sel Hijau(1,2) memiliki \textit{support} pada Kuning(6,3) dan tetap valid.}
	\label{fig:bt-step24}
\end{figure}

Sebaliknya, Gambar~\ref{fig:bt-step25} menunjukkan pemrosesan sel Hijau(5,2) dan Hijau(5,3). Sel Hijau(5,2) berkonflik dengan Kuning(6,2) karena berada pada kolom yang sama, sekaligus berkonflik dengan Kuning(6,3) karena bertetangga secara diagonal langsung. Demikian pula, sel Hijau(5,3) berkonflik dengan Kuning(6,2) secara diagonal dan dengan Kuning(6,3) pada kolom yang sama. Karena tidak ada satu pun sel pada domain warna kuning yang dapat menjadi \textit{support}, kedua sel hijau tersebut dipangkas dari domain dan ditandai dengan lingkaran ungu pada papan.

\begin{figure}[!htbp]
	\centering
	\includegraphics[width=0.27\linewidth]{./bt/25.png}
	\caption{Pemrosesan busur Hijau$\rightarrow$Kuning: sel Hijau(5,2) dan Hijau(5,3) tidak memiliki \textit{support} dan dipangkas.}
	\label{fig:bt-step25}
\end{figure}

Propagasi serupa dilakukan untuk seluruh busur dalam antrean, dengan busur-busur baru ditambahkan setiap kali sebuah domain mengalami pemangkasan. Kondisi papan setelah seluruh propagasi AC-3 selesai diperlihatkan pada Gambar~\ref{fig:bt-step26}. Tanda X merah menunjukkan sel yang dipangkas oleh \textit{forward checking}, sedangkan tanda X dengan lingkaran ungu menunjukkan sel tambahan yang dipangkas oleh AC-3. Penggabungan kedua mekanisme ini menghasilkan reduksi ruang pencarian yang lebih agresif dibandingkan menggunakan salah satunya saja.

\begin{figure}[!htbp]
	\centering
	\includegraphics[width=0.27\linewidth]{./bt/26.png}
	\caption{Kondisi papan setelah propagasi AC-3 selesai.}
	\label{fig:bt-step26}
\end{figure}

\subsubsection{Penempatan Berikutnya dan Deteksi \textit{Domain Wipeout}}

Setelah pemangkasan selesai dan tidak ada domain yang kosong, algoritma melanjutkan pencarian dengan memanggil dirinya sendiri secara rekursif untuk variabel warna berikutnya. Misalkan algoritma menempatkan menteri pada salah satu sel dari domain warna oranye sebagaimana ditunjukkan pada Gambar~\ref{fig:bt-step27}.

\begin{figure}[!htbp]
	\centering
	\includegraphics[width=0.27\linewidth]{./bt/27.png}
	\caption{Penempatan menteri kedua pada salah satu sel dari domain warna oranye.}
	\label{fig:bt-step27}
\end{figure}

Setelah penempatan ini, \textit{forward checking} dan AC-3 dijalankan kembali. Namun, kombinasi penempatan ini ternyata menyebabkan seluruh sel pada domain warna kuning menjadi tidak valid, sebagaimana diperlihatkan pada Gambar~\ref{fig:bt-step28}. Kondisi ini merupakan \textit{domain wipeout}, di mana tidak ada lagi nilai yang dapat ditugaskan kepada warna kuning, sehingga \textit{backtrack} diperlukan.

\begin{figure}[!htbp]
	\centering
	\includegraphics[width=0.27\linewidth]{./bt/28.png}
	\caption{\textit{Domain wipeout}: seluruh sel pada domain warna kuning menjadi tidak valid.}
	\label{fig:bt-step28}
\end{figure}

\subsubsection{\textit{Backtrack} dan Pemulihan Domain}

Sebagai respons terhadap kegagalan, algoritma melakukan \textit{backtrack} dengan membatalkan penempatan menteri oranye yang baru saja dilakukan, sebagaimana diperlihatkan pada Gambar~\ref{fig:bt-step29}. Seluruh pemangkasan yang dilakukan akibat penempatan tersebut dipulihkan, mengembalikan domain seluruh variabel ke kondisi tepat sebelum penempatan oranye dilakukan.

\begin{figure}[!htbp]
	\centering
	\includegraphics[width=0.27\linewidth]{./bt/29.png}
	\caption{\textit{Backtrack}: pembatalan penempatan menteri oranye yang menyebabkan kegagalan.}
	\label{fig:bt-step29}
\end{figure}

Algoritma kemudian mencoba sel berikutnya dari domain warna oranye yang masih valid, sebagaimana diperlihatkan pada Gambar~\ref{fig:bt-step30}. Pemangkasan dijalankan kembali pada konfigurasi baru ini, dan apabila tidak terjadi \textit{domain wipeout}, pencarian dilanjutkan ke variabel warna berikutnya.

\begin{figure}[!htbp]
	\centering
	\includegraphics[width=0.27\linewidth]{./bt/30.png}
	\caption{Penempatan alternatif untuk warna oranye setelah \textit{backtrack}.}
	\label{fig:bt-step30}
\end{figure}

\subsubsection{Solusi Akhir}

Proses penempatan, pemangkasan, dan \textit{backtrack} berulang hingga seluruh tujuh variabel warna mendapat penugasan yang konsisten. Solusi akhir untuk papan 7$\times$7 diperlihatkan pada Gambar~\ref{fig:bt-step31}, di mana setiap sektor warna memiliki tepat satu menteri yang ditempatkan, tidak ada dua menteri yang berbagi baris atau kolom yang sama, dan tidak ada dua menteri yang bertetangga secara langsung.

\begin{figure}[!htbp]
	\centering
	\includegraphics[width=0.4\linewidth]{./bt/31.png}
	\caption{Solusi akhir papan 7$\times$7 yang memenuhi seluruh kendala \textit{Colored Queens}.}
	\label{fig:bt-step31}
\end{figure}

Studi kasus ini mengilustrasikan bagaimana sinergi antara \textit{forward checking} dan AC-3 secara signifikan mengurangi ruang pencarian. \textit{Forward checking} menangani konflik langsung dengan penugasan terbaru, sementara AC-3 menelusuri konsekuensi tidak langsung dari penugasan tersebut terhadap konsistensi seluruh pasangan variabel yang belum ditugaskan. Mekanisme \textit{backtrack} memastikan bahwa kegagalan pada satu cabang pencarian dapat ditangani secara efisien dengan memulihkan domain ke kondisi sebelumnya, sehingga algoritma dapat mengeksplorasi cabang alternatif tanpa perlu menghitung ulang seluruh struktur domain.


\section{Analisis Algoritma \textit{Particle Swarm Optimization}}
\label{sec:analisis-pso}

Bagian ini menjelaskan analisis algoritmik dari pendekatan kedua yang digunakan untuk menyelesaikan permasalahan \textit{Colored Queens}, yaitu \textit{Particle Swarm Optimization} (PSO) dalam varian diskrit. Berbeda dengan \textit{backtracking} yang menjamin penemuan solusi melalui pencarian sistematis, PSO merupakan algoritma stokastik berbasis populasi yang mencari solusi melalui pergerakan kolektif partikel-partikel pada ruang pencarian. Landasan teori PSO klasik dan adaptasinya untuk ruang pencarian diskrit telah dibahas pada Bagian~\ref{subsubsec:pso}. Analisis disajikan melalui dua diagram alur, yaitu mekanisme pembaruan kecepatan dan posisi setiap partikel, kemudian alur utama algoritma yang menjalankan pembaruan tersebut secara iteratif hingga salah satu kondisi terminasi tercapai. Diagram pembaruan partikel disajikan terlebih dahulu agar pembaca dapat memahami unit komputasi inti dari PSO sebelum melihat keseluruhan alur. Rincian implementasi masing-masing komponen dibahas pada Bab~\ref{chap:perancangan}.

\subsection{Fungsi \textit{Fitness}}
\label{subsec:fungsi-fitness-pso}

Berbeda dengan algoritma \textit{backtracking} yang menjamin validitas solusi melalui pemeriksaan kendala secara eksplisit selama proses pencarian, PSO sebagai algoritma berbasis populasi tidak menggunakan kendala secara langsung untuk membatasi pergerakan partikel. Sebagai gantinya, kualitas suatu posisi partikel diukur secara numerik melalui sebuah fungsi \textit{fitness}, yaitu fungsi yang memetakan setiap posisi partikel ke sebuah nilai real yang mencerminkan seberapa dekat posisi tersebut dengan solusi valid. Pendekatan ini sejalan dengan transformasi umum dari \textit{Constraint Satisfaction Problem} (CSP) menjadi \textit{Constraint Optimization Problem} (COP) sebagaimana telah dibahas pada Bagian~\ref{subsubsec:metaheuristic-konsep}: dengan reformulasi ini, permasalahan ``mencari penugasan yang memenuhi seluruh kendala'' diubah menjadi permasalahan ``meminimalkan jumlah pelanggaran kendala,'' yang dapat ditangani oleh \textit{metaheuristic} berbasis optimasi seperti PSO. Posisi partikel yang merepresentasikan solusi valid bagi permasalahan \textit{Colored Queens} berkorespondensi dengan nilai \textit{fitness} sama dengan nol, dan algoritma berhenti ketika kondisi tersebut tercapai.

Fungsi \textit{fitness} pada penelitian ini dirancang untuk mengukur dua jenis pelanggaran kendala yang dapat terjadi pada konfigurasi penempatan menteri, yang masing-masing bersesuaian dengan kendala yang telah didefinisikan pada Bagian~\ref{subsec:definisi-constraint}. Pelanggaran pertama, yang disebut \textit{attacking violation}, menghitung jumlah pasangan menteri yang ditempatkan pada baris atau kolom yang sama. Pelanggaran jenis ini menggabungkan kendala baris dan kendala kolom yang secara struktural setara: keduanya melarang dua menteri berbagi sebuah sumbu linear pada papan. Pelanggaran kedua, yang disebut \textit{adjacency violation}, menghitung jumlah pasangan menteri yang menempati sel-sel yang bertetangga secara langsung dalam delapan arah. Kedua jenis pelanggaran ini diperhitungkan secara terpisah karena merepresentasikan dimensi konflik yang berbeda secara semantik: \textit{attacking} mengukur konflik berbasis sumbu (\textit{axis-based conflict}), sedangkan \textit{adjacency} mengukur konflik berbasis kedekatan posisi (\textit{proximity-based conflict}). Kendala sektor warna tidak perlu diperhitungkan dalam fungsi \textit{fitness} karena telah dipenuhi secara implisit oleh representasi posisi partikel, di mana setiap indeks pada vektor posisi merepresentasikan pilihan sel dari domain satu sektor tertentu.

Kedua komponen pelanggaran tersebut kemudian dikombinasikan menjadi sebuah nilai \textit{fitness} tunggal melalui penjumlahan berbobot, sebagaimana ditunjukkan pada Persamaan~\ref{eq:fitness-pso}.

\begin{equation}
	f(\vec{x}) = w_1 \cdot V_{\textit{attacking}}(\vec{x}) + w_2 \cdot V_{\textit{adjacency}}(\vec{x})
	\label{eq:fitness-pso}
\end{equation}

di mana $\vec{x}$ adalah posisi partikel, $V_{\textit{attacking}}(\vec{x})$ adalah jumlah pasangan menteri yang melanggar kendala \textit{attacking}, $V_{\textit{adjacency}}(\vec{x})$ adalah jumlah pasangan menteri yang melanggar kendala \textit{adjacency}, serta $w_1$ dan $w_2$ adalah bobot positif yang menentukan kontribusi relatif setiap jenis pelanggaran terhadap nilai \textit{fitness} keseluruhan. Pemilihan bobot mencerminkan prioritas relatif dalam penyelesaian: nilai $w_1$ yang lebih besar mendorong algoritma untuk memprioritaskan pengurangan pelanggaran \textit{attacking}, sedangkan nilai $w_2$ yang lebih besar mendorong prioritas pada pengurangan pelanggaran \textit{adjacency}. Konfigurasi \textit{baseline} pada tugas akhir ini menggunakan $w_1 = w_2 = 1$, yang memberikan bobot setara pada kedua jenis pelanggaran dan menjadikan fungsi \textit{fitness} setara dengan jumlah total pelanggaran kendala pada konfigurasi tersebut.

Penjumlahan berbobot dipilih sebagai mekanisme agregasi karena kesederhanaan interpretasinya dan kompatibilitasnya dengan kerangka kerja optimasi PSO yang bersifat skalar. Fungsi \textit{fitness} sebagai keluaran skalar memungkinkan perbandingan antarposisi partikel secara langsung melalui operator perbandingan numerik, yang diperlukan untuk pembaruan \textit{pBest} dan \textit{nBest} pada setiap iterasi. Alternatif lain seperti agregasi multiplikatif atau formulasi vektorial akan menghasilkan permasalahan multi-objektif yang memerlukan mekanisme penanganan tambahan, seperti konsep dominasi Pareto, yang menambah kompleksitas algoritma tanpa memberikan keuntungan signifikan pada konteks \textit{Colored Queens} di mana kedua jenis pelanggaran sama-sama harus dieliminasi secara absolut pada solusi valid.

Perlu dicatat bahwa kedua komponen pelanggaran tidak bersifat saling-eksklusif. Sepasang menteri yang berada pada baris atau kolom yang sama dan secara bersamaan menempati sel-sel yang bertetangga akan dihitung sebagai pelanggaran pada kedua komponen secara independen. Sebagai contoh, dua menteri yang ditempatkan pada sel $(r, c)$ dan $(r, c+1)$ melanggar kendala baris (keduanya pada baris $r$) sekaligus kendala \textit{adjacency} (keduanya bertetangga secara horizontal). Penghitungan ganda ini bukan suatu redundansi, melainkan keputusan desain yang disengaja: pasangan yang melanggar dua kendala sekaligus merepresentasikan konfigurasi yang lebih jauh dari solusi valid dibandingkan pasangan yang hanya melanggar satu kendala, sehingga sewajarnya menerima penalti yang lebih tinggi pada nilai \textit{fitness}. Mekanisme ini menghasilkan gradien yang lebih curam pada konfigurasi yang sangat melanggar, yang membantu mengarahkan partikel keluar dari area pencarian yang tidak menjanjikan dengan lebih cepat.

\subsection{Pembaruan Kecepatan dan Posisi Partikel}
\label{subsec:flow-particle-update}

Mekanisme pembaruan kecepatan dan posisi merupakan inti komputasi PSO yang dijalankan pada setiap iterasi terhadap seluruh partikel dalam populasi. Pada varian diskrit yang digunakan dalam tugas akhir ini, kecepatan tidak diinterpretasikan sebagai perubahan posisi numerik, melainkan sebagai probabilitas adopsi posisi terbaik tetangga (\textit{neighborhood best} atau \textit{nBest}) pada setiap dimensi. Pendekatan ini mempertahankan kerangka kerja PSO klasik sekaligus mengakomodasi sifat diskrit dari ruang pencarian \textit{Colored Queens}. Alur pembaruan partikel diperlihatkan pada Gambar~\ref{fig:flow-particle-update}.

\begin{figure}[!htbp]
	\centering
	\includegraphics[width=0.95\linewidth]{flowparticle.png}
	\caption{Diagram alur pembaruan kecepatan dan posisi partikel PSO diskrit.}
	\label{fig:flow-particle-update}
\end{figure}

\subsubsection{Iterasi atas \textit{Neighborhood}}

Prosedur dimulai dengan mengiterasi seluruh kelompok \textit{neighborhood} yang telah dibentuk pada tahap inisialisasi populasi. Setiap \textit{neighborhood} adalah himpunan partikel yang berbagi informasi mengenai posisi terbaik kolektifnya (\textit{nBest}), sehingga partikel-partikel dalam satu kelompok bergerak dipandu oleh referensi yang sama. Penggunaan topologi \textit{neighborhood} dipilih atas pertimbangan yang dijabarkan pada Bagian~\ref{subsubsec:pso}. Apabila seluruh \textit{neighborhood} telah diproses, prosedur berakhir.

\subsubsection{Iterasi atas Partikel dalam \textit{Neighborhood}}

Untuk setiap \textit{neighborhood} yang diambil, prosedur mengiterasi seluruh partikel yang menjadi anggotanya. Setiap partikel akan melalui rangkaian langkah pembaruan secara independen, tetapi semuanya merujuk pada \textit{nBest} yang sama untuk kelompok tersebut. Apabila seluruh partikel dalam \textit{neighborhood} saat ini telah diperbarui, alur kembali ke iterasi \textit{neighborhood} berikutnya.

\subsubsection{Penghitungan Kecepatan Baru}
\label{subsubsec:update-velocity}
Kecepatan baru setiap partikel dihitung sebagai kombinasi linear dari tiga komponen, yaitu komponen inersia yang mempertahankan arah pergerakan sebelumnya, komponen kognitif yang menarik partikel menuju posisi terbaik historisnya (\textit{personal best} atau \textit{pBest}), dan komponen sosial yang menarik partikel menuju \textit{nBest}. Untuk setiap partikel, kecepatan baru pada dimensi ke-$j$ dihitung sebagaimana ditunjukkan pada Persamaan~\ref{eq:pso-velocity-update}.

\begin{equation}
	v_j^{(t+1)} = w \cdot v_j^{(t)} + c_1 \cdot r_1 \cdot \delta_{pBest,j} + c_2 \cdot r_2 \cdot \delta_{nBest,j}
	\label{eq:pso-velocity-update}
\end{equation}

di mana $w$ adalah koefisien inersia, $c_1$ dan $c_2$ adalah koefisien akselerasi kognitif dan sosial, $r_1$ dan $r_2$ adalah bilangan acak seragam pada interval $[0, 1)$ yang dibangkitkan secara independen untuk setiap dimensi, dan $\delta_{pBest,j}$ serta $\delta_{nBest,j}$ adalah indikator perbedaan posisi yang didefinisikan pada Persamaan~\ref{eq:pso-delta}.

\begin{equation}
	\delta_{pBest,j} = \begin{cases} 1 & \text{jika } x_j^{(t)} \neq pBest_j \\ 0 & \text{jika } x_j^{(t)} = pBest_j \end{cases}
	\qquad
	\delta_{nBest,j} = \begin{cases} 1 & \text{jika } x_j^{(t)} \neq nBest_j \\ 0 & \text{jika } x_j^{(t)} = nBest_j \end{cases}
	\label{eq:pso-delta}
\end{equation}

Penggunaan indikator biner $\delta$ sebagai pengganti selisih aritmatika merupakan adaptasi kunci dari PSO kontinu ke ruang diskrit. Pada PSO kontinu, komponen kognitif dan sosial dihitung melalui selisih vektor posisi yang menghasilkan arah perpindahan dalam ruang Euclidean. Pada ruang diskrit berindeks, selisih aritmatika antara dua indeks sel tidak memiliki makna geometris: sel dengan indeks tertentu dan sel dengan indeks lain dalam domain suatu warna tidak memiliki hubungan spasial yang ditentukan oleh selisih indeksnya. Oleh sebab itu, informasi yang relevan hanyalah apakah posisi saat ini berbeda dari posisi referensi (\textit{pBest} atau \textit{nBest}), bukan seberapa jauh perbedaannya.

Setelah perhitungan, nilai kecepatan di-\textit{clamp} pada interval $[0, 1]$ untuk mempertahankan interpretasinya sebagai probabilitas, sebagaimana ditunjukkan pada Persamaan~\ref{eq:pso-clamp}.

\begin{equation}
	v_j^{(t+1)} = \max(0,\; \min(1,\; v_j^{(t+1)}))
	\label{eq:pso-clamp}
\end{equation}

\subsubsection{Pembaruan Posisi Partikel}
\label{subsubsec:update-position}
Posisi baru partikel diturunkan dari kecepatan dengan interpretasi probabilistik. Untuk setiap dimensi $j$, sebuah bilangan acak $r \in [0, 1)$ dibangkitkan, kemudian dibandingkan dengan nilai kecepatan pada dimensi tersebut. Apabila $r < v_j$, posisi partikel pada dimensi tersebut diadopsi dari \textit{nBest}. Apabila $r \geq v_j$, posisi tidak berubah. Secara formal, mekanisme pembaruan posisi ini dinyatakan pada Persamaan~\ref{eq:pso-position-update}.

\begin{equation}
	x_j^{(t+1)} = \begin{cases} nBest_j & \text{jika } r < v_j^{(t+1)} \\ x_j^{(t)} & \text{jika } r \geq v_j^{(t+1)} \end{cases}
	\label{eq:pso-position-update}
\end{equation}

Mekanisme ini secara sengaja hanya menggunakan posisi \textit{nBest} sebagai target perpindahan, meskipun \textit{pBest} turut memengaruhi probabilitas perpindahan melalui formula kecepatan. Pemisahan peran ini merupakan keputusan desain yang didasarkan pada karakteristik ruang pencarian diskrit. Pada PSO kontinu, posisi baru merupakan hasil penjumlahan vektor yang secara alami menginterpolasi antara \textit{pBest} dan \textit{nBest}. Pada ruang diskrit, interpolasi semacam ini tidak terdefinisi: sebuah dimensi hanya dapat mengambil satu nilai indeks pada satu waktu, sehingga partikel harus memilih salah satu target, bukan menggabungkan keduanya. Apabila pembaruan posisi juga mempertimbangkan \textit{pBest} sebagai target alternatif, setiap dimensi akan menghadapi dua arah perpindahan yang dapat saling bertentangan, yang dalam ruang diskrit menghasilkan pergantian indeks yang tidak koheren alih-alih konvergensi bertahap. Dengan menjadikan \textit{nBest} sebagai satu-satunya target perpindahan, mekanisme pembaruan posisi menjadi lebih terarah, sementara pengaruh pengalaman individual partikel tetap terkodifikasi dalam besaran probabilitas melalui komponen $c_1 \cdot r_1 \cdot \delta_{pBest,j}$ pada formula kecepatan.

\subsubsection{Evaluasi \textit{Fitness} dan Pembaruan \textit{pBest} serta \textit{nBest}}

Setelah posisi partikel diperbarui, nilai \textit{fitness} dievaluasi ulang berdasarkan fungsi \textit{fitness} yang menghitung jumlah pelanggaran kendala sesuai definisi pada Bagian~\ref{subsec:definisi-constraint} dan formulasi pada Bagian~\ref{subsec:fungsi-fitness-pso}. Nilai \textit{fitness} ini kemudian dibandingkan dengan \textit{fitness} dari \textit{pBest} partikel: apabila lebih baik, \textit{pBest} diperbarui dengan posisi saat ini. Selanjutnya, \textit{fitness} juga dibandingkan dengan \textit{fitness} dari \textit{nBest} \textit{neighborhood}: apabila lebih baik, \textit{nBest} diperbarui agar partikel-partikel lain dalam kelompok yang sama dapat mengikutinya pada iterasi berikutnya.

\subsection{Alur Utama Algoritma PSO}
\label{subsec:flow-pso-main}

Alur utama algoritma PSO mengorkestrasi keseluruhan proses optimasi, mulai dari inisialisasi populasi hingga terminasi. Pada setiap iterasi, mekanisme pembaruan partikel yang telah dijabarkan pada Bagian~\ref{subsec:flow-particle-update} dijalankan terhadap seluruh partikel, diikuti oleh evaluasi terhadap tiga kondisi terminasi: ditemukannya solusi valid, tercapainya batas iterasi maksimum, atau terdeteksinya stagnasi. Keseluruhan alur diperlihatkan pada Gambar~\ref{fig:flow-pso-main}.

\begin{figure}[!htbp]
	\centering
	\includegraphics[width=0.95\linewidth]{flowpso.png}
	\caption{Diagram alur utama algoritma PSO diskrit.}
	\label{fig:flow-pso-main}
\end{figure}

\subsubsection{Inisialisasi Populasi dan Variabel Pelacakan}

Algoritma dimulai dengan inisialisasi populasi, yaitu pembangkitan sejumlah partikel dengan posisi acak pada ruang pencarian dan kecepatan acak pada interval valid. Selanjutnya, partikel-partikel tersebut dipartisi ke dalam beberapa \textit{neighborhood} dengan ukuran yang relatif seimbang, dan nilai \textit{fitness} awal seluruh partikel dievaluasi. Setelah inisialisasi populasi selesai, algoritma menetapkan pencacah iterasi pada nilai satu, pencacah stagnasi pada nilai nol, dan menyimpan \textit{fitness} terendah di antara seluruh \textit{nBest} sebagai nilai \textit{fitness} global terbaik (\textit{globalBest}) yang akan digunakan sebagai acuan deteksi stagnasi pada iterasi-iterasi berikutnya.

\subsubsection{Pengecekan Batas Iterasi}

Pada awal setiap iterasi, algoritma memeriksa apakah pencacah iterasi telah melampaui batas iterasi maksimum yang telah ditetapkan sebagai parameter. Apabila ya, pencarian dihentikan dan algoritma melaporkan partikel dengan \textit{fitness} terendah di antara seluruh \textit{nBest} sebagai solusi terbaik yang ditemukan, meskipun solusi tersebut belum tentu valid. Apabila batas belum tercapai, algoritma melanjutkan ke tahap pembaruan partikel.

\subsubsection{Pembaruan Seluruh Partikel}

Algoritma menjalankan dua tahap pembaruan secara berurutan terhadap seluruh partikel dalam populasi. Tahap pertama memperbarui kecepatan seluruh partikel berdasarkan \textit{pBest} masing-masing dan \textit{nBest} dari \textit{neighborhood}-nya. Tahap kedua memperbarui posisi partikel berdasarkan kecepatan yang baru dihitung, mengevaluasi ulang \textit{fitness}, serta memperbarui \textit{pBest} dan \textit{nBest} apabila ditemukan posisi yang lebih baik. Kedua tahap ini secara kolektif merupakan eksekusi penuh dari prosedur pembaruan partikel yang telah dijabarkan pada Bagian~\ref{subsec:flow-particle-update}.

\subsubsection{Pengecekan Solusi Valid}

Setelah pembaruan selesai, algoritma memeriksa apakah salah satu \textit{nBest} memiliki nilai \textit{fitness} sama dengan nol. Nilai \textit{fitness} nol menandakan tidak adanya pelanggaran kendala pada posisi tersebut, yang berarti partikel tersebut merepresentasikan solusi valid bagi permasalahan \textit{Colored Queens}. Pemeriksaan dilakukan hanya terhadap \textit{nBest} dari setiap \textit{neighborhood}, bukan terhadap seluruh partikel, karena partikel dengan \textit{fitness} terendah dalam setiap kelompok pasti merupakan \textit{nBest} saat ini. Apabila solusi valid ditemukan, algoritma berhenti dan melaporkan solusi tersebut.

\subsubsection{Pembaruan \textit{Fitness} Global dan Pencacah Stagnasi}

Apabila solusi valid belum ditemukan, algoritma menghitung \textit{fitness} terendah di antara seluruh \textit{nBest} saat ini (\textit{fitnessCur}) dan membandingkannya dengan \textit{globalBest} yang tersimpan. Apabila \textit{fitnessCur} lebih rendah, hal ini menandakan bahwa populasi mengalami perbaikan, sehingga \textit{globalBest} diperbarui dan pencacah stagnasi direset menjadi nol. Apabila tidak terjadi perbaikan, pencacah stagnasi ditambahkan satu untuk menandai berlalunya satu iterasi tanpa kemajuan.

\subsubsection{Pengecekan Stagnasi dan Lanjutan Iterasi}

Setelah pencacah stagnasi diperbarui, algoritma memeriksa apakah nilainya telah melampaui batas stagnasi maksimum. Mekanisme ini menghentikan pencarian lebih awal pada kasus di mana populasi telah terjebak pada \textit{local optimum} dan eksplorasi lebih lanjut tidak produktif. Justifikasi keberadaan dan pemilihan nilai batas stagnasi dijabarkan pada Bagian~\ref{subsec:stagnasi-terminasi}. Apabila batas tidak terlampaui, pencacah iterasi ditambahkan satu dan alur kembali ke pengecekan batas iterasi, melanjutkan siklus optimasi pada iterasi berikutnya.

\subsection{Studi Kasus: Pencarian Solusi dengan PSO pada Papan 7$\times$7}
\label{subsec:studi-kasus-pso}

Sebagai pembanding terhadap pendekatan \textit{backtracking} yang telah dijabarkan pada Bagian~\ref{subsec:studi-kasus-bt}, bagian ini menyajikan studi kasus penerapan PSO pada papan 7$\times$7 yang sama. Untuk menjaga kesederhanaan ilustrasi, contoh ini menggunakan populasi minimal yang terdiri atas dua partikel yang dikelompokkan dalam satu \textit{neighborhood}. Pada implementasi nyata, jumlah partikel dan \textit{neighborhood} jauh lebih besar untuk meningkatkan keragaman eksplorasi.

\subsubsection{Inisialisasi Populasi}

Pada tahap inisialisasi, dua partikel dibangkitkan dengan posisi awal acak. Setiap posisi partikel direpresentasikan sebagai vektor berukuran tujuh, di mana setiap index pada vektor berkaitan dengan satu sektor warna. Nilai pada index ke-$i$ menyatakan posisi sel terpilih dari domain warna ke-$i$. Posisi awal kedua partikel ditetapkan sebagai berikut:

\begin{itemize}
	\item Partikel 1: $\{3,\, 1,\, 12,\, 8,\, 0,\, 2,\, 1\}$
	\item Partikel 2: $\{0,\, 3,\, 9,\, 2,\, 0,\, 3,\, 1\}$
\end{itemize}

Selain posisi, setiap partikel juga memiliki vektor kecepatan awal dengan nilai acak pada interval $[0, 1)$ untuk setiap index:

\begin{itemize}
	\item Partikel 1: $v_i = [0.12,\, 0.85,\, 0.43,\, 0.67,\, 0.09,\, 0.91,\, 0.274]$
	\item Partikel 2: $v_i = [0.76,\, 0.21,\, 0.58,\, 0.04,\, 0.69,\, 0.47,\, 0.83]$
\end{itemize}

\subsubsection{Perhitungan \textit{Fitness} Awal dan Penetapan \textit{nBest}}

Nilai \textit{fitness} awal setiap partikel dievaluasi berdasarkan fungsi \textit{fitness} yang telah dijabarkan pada Bagian~\ref{subsec:fungsi-fitness-pso}. Dalam contoh ini, hasil evaluasi adalah:

\begin{itemize}
	\item Partikel 1: $f = 14$
	\item Partikel 2: $f = 20$
\end{itemize}

Karena partikel 1 memiliki nilai \textit{fitness} terendah, partikel 1 ditetapkan sebagai \textit{neighborhood best} (\textit{nBest}) dengan nilai $14$. Penetapan ini akan digunakan sebagai acuan untuk pembaruan kecepatan dan posisi pada iterasi berikutnya. Selain itu, posisi awal masing-masing partikel juga ditetapkan sebagai \textit{personal best} (\textit{pBest}) mereka.

\subsubsection{Pembaruan Kecepatan}

Pada iterasi pertama, kecepatan setiap partikel diperbarui menggunakan Persamaan~\ref{eq:pso-velocity-update}. Sebagai contoh ilustratif, misalkan parameter PSO ditetapkan sebagai $\omega = 0.7$, $c_1 = 0.5$, $c_2 = 0.5$, dan bilangan acak yang dibangkitkan adalah $r_1 = 0.2$ serta $r_2 = 0.27$. Untuk index ke-0 dari partikel 2, perhitungan kecepatan baru ditunjukkan pada Persamaan~\ref{eq:contoh-velocity-pso} dan menghasilkan nilai $v_{i+1}[0] = 0.682$.

\begin{equation}
	\begin{aligned}
		v_{i+1}[0] &= \underbrace{(\omega \cdot v_i[0])}_{\text{inersia}} + \underbrace{(r_1 \cdot c_1 \cdot \delta_{pBest,0})}_{\text{kognitif}} + \underbrace{(r_2 \cdot c_2 \cdot \delta_{nBest,0})}_{\text{sosial}} \\
		&= (0.7 \cdot 0.76) + (0.2 \cdot 0.5 \cdot 0) + (0.27 \cdot 0.5 \cdot 1) \\
		&= 0.532 + 0 + 0.15 = 0.682
	\end{aligned}
	\label{eq:contoh-velocity-pso}
\end{equation}

Nilai $\delta_{pBest,0} = 0$ pada Persamaan~\ref{eq:contoh-velocity-pso} muncul karena nilai partikel 2 pada index ke-0 sama dengan nilai \textit{pBest}-nya pada index yang sama (keduanya bernilai $0$), sesuai dengan definisi $\delta$ pada Persamaan~\ref{eq:pso-delta}. Sebaliknya, $\delta_{nBest,0} = 1$ karena nilai partikel 2 pada index ke-0 (bernilai $0$) berbeda dengan nilai \textit{nBest} pada index ke-0 (bernilai $3$). Perhitungan serupa dilakukan untuk seluruh index, menghasilkan vektor kecepatan baru pada Persamaan~\ref{eq:contoh-velocity-vector}.

\begin{equation}
	v_{i+1} = [0.682,\, 0.497,\, 0.606,\, 0.2728,\, 0.483,\, 0.579,\, 0.581]
	\label{eq:contoh-velocity-vector}
\end{equation}

\subsubsection{Pembaruan Posisi}

Selanjutnya, posisi partikel diperbarui berdasarkan kecepatan baru menggunakan Persamaan~\ref{eq:pso-position-update}. Untuk setiap index, bilangan acak $r \in [0, 1)$ dibangkitkan, dan apabila $r < v_{i+1}[j]$, nilai pada index tersebut diadopsi dari \textit{nBest}; apabila tidak, nilai tetap.

Sebagai contoh, untuk index ke-0 partikel 2, dibangkitkan bilangan acak $r = 0.25$. Karena $r < v_{i+1}[0]$ (yaitu $0.25 < 0.682$), maka nilai pada index ke-0 mengadopsi nilai dari \textit{nBest}, sehingga $x_{i+1}[0] = 3$ (menggantikan nilai sebelumnya yang bernilai $0$). Perhitungan serupa dilakukan untuk seluruh index, menghasilkan posisi baru partikel 2 pada Persamaan~\ref{eq:contoh-posisi-baru}.

\begin{equation}
	x_{i+1} = [3,\, 3,\, 12,\, 2,\, 0,\, 2,\, 1]
	\label{eq:contoh-posisi-baru}
\end{equation}

\subsubsection{Evaluasi \textit{Fitness} Partikel Baru}

Setelah posisi diperbarui, nilai \textit{fitness} partikel 2 dievaluasi ulang menggunakan Persamaan~\ref{eq:fitness-pso}. Untuk contoh ini, ditetapkan $w_1 = 1$ dan $w_2 = 2$. Posisi baru partikel 2 pada Persamaan~\ref{eq:contoh-posisi-baru} memetakan menteri ke sel-sel tertentu pada papan, dan pelanggaran kendala kemudian dihitung berdasarkan dua komponen, yaitu $V_{\textit{attacking}}$ dan $V_{\textit{adjacency}}$.

Dengan kelima pasangan tersebut, nilai $V_{\textit{attacking}}$ adalah $5$. Nilai \textit{fitness} akhir partikel 2 setelah iterasi pertama dihitung berdasarkan Persamaan~\ref{eq:fitness-pso} dan menghasilkan nilai $11$, sebagaimana ditunjukkan pada Persamaan~\ref{eq:contoh-fitness-1}.

\begin{equation}
	f(\vec{x}) = w_1 \cdot V_{\textit{attacking}} + w_2 \cdot V_{\textit{adjacency}} = (1 \cdot 5) + (2 \cdot 3) = 5 + 6 = 11
	\label{eq:contoh-fitness-1}
\end{equation}

\textbf{Pelanggaran \textit{adjacency}.} Komponen ini menghitung pasangan menteri yang bertetangga secara langsung dalam delapan arah. Pada konfigurasi partikel 2, terdapat tiga pasangan \textit{adjacency}. Pasangan pertama, sebagaimana disorot pada Gambar~\ref{fig:pso-step39}, terdiri atas menteri pada (3,1) dan menteri pada (4,1) yang bertetangga secara vertikal.

\begin{figure}[!htbp]
	\centering
	\includegraphics[width=0.27\linewidth]{./pso/39.png}
	\caption{Pasangan \textit{adjacency} pertama: menteri pada (3,1) dan (4,1) bertetangga vertikal.}
	\label{fig:pso-step39}
\end{figure}

Pasangan kedua, sebagaimana disorot pada Gambar~\ref{fig:pso-step40}, terdiri atas menteri pada (4,1) dan menteri pada (4,2) yang bertetangga secara horizontal.

\begin{figure}[!htbp]
	\centering
	\includegraphics[width=0.27\linewidth]{./pso/40.png}
	\caption{Pasangan \textit{adjacency} kedua: menteri pada (4,1) dan (4,2) bertetangga horizontal.}
	\label{fig:pso-step40}
\end{figure}

Pasangan ketiga, sebagaimana disorot pada Gambar~\ref{fig:pso-step41}, terdiri atas menteri pada (3,1) dan menteri pada (4,2) yang bertetangga secara diagonal.

\begin{figure}[!htbp]
	\centering
	\includegraphics[width=0.27\linewidth]{./pso/41.png}
	\caption{Pasangan \textit{adjacency} ketiga: menteri pada (3,1) dan (4,2) bertetangga diagonal.}
	\label{fig:pso-step41}
\end{figure}

Dengan demikian, nilai \textit{adjacencyViolation} adalah $3$..

\textbf{Pelanggaran \textit{attacking}.} Komponen ini menghitung pasangan menteri yang berada pada baris atau kolom yang sama. Pada konfigurasi partikel 2, Gambar~\ref{fig:pso-step42} menyoroti pasangan pertama, yaitu menteri pada (4,2) dan menteri pada (6,2) yang berada pada kolom 2. 

\begin{figure}[!htbp]
	\centering
	\includegraphics[width=0.27\linewidth]{./pso/42.png}
	\caption{Pasangan \textit{attacking} pertama: menteri pada (4,2) dan (6,2) berada pada kolom yang sama.}
	\label{fig:pso-step42}
\end{figure}

Pasangan kedua, sebagaimana disorot pada Gambar~\ref{fig:pso-step43}, terdiri atas menteri pada (4,1) dan menteri pada (4,2) yang berada pada baris 4. Perlu dicatat bahwa pasangan ini juga merupakan pelanggaran \textit{adjacency} sebagaimana telah disorot pada Gambar~\ref{fig:pso-step40}. Pasangan tersebut tetap dihitung pada kedua komponen secara independen karena keduanya merepresentasikan jenis kendala yang berbeda: \textit{attacking} mengukur konflik berbasis baris dan kolom, sedangkan \textit{adjacency} mengukur kedekatan posisi.

\begin{figure}[!htbp]
	\centering
	\includegraphics[width=0.27\linewidth]{./pso/46.png}
	\caption{Pasangan \textit{attacking} kedua: menteri pada (4,1) dan (4,2) berada pada baris yang sama. Pasangan ini juga merupakan pelanggaran \textit{adjacency}.}
	\label{fig:pso-step43}
\end{figure}

Pasangan ketiga, sebagaimana disorot pada Gambar~\ref{fig:pso-step44}, terdiri atas menteri pada (4,1) dan menteri pada (4,4) yang berada pada baris 4. Berbeda dengan pasangan sebelumnya, pasangan ini tidak bertetangga langsung sehingga hanya tergolong sebagai pelanggaran \textit{attacking}.

\begin{figure}[!htbp]
	\centering
	\includegraphics[width=0.27\linewidth]{./pso/44.png}
	\caption{Pasangan \textit{attacking} ketiga: menteri pada (4,1) dan (4,4) berada pada baris yang sama.}
	\label{fig:pso-step44}
\end{figure}

Pasangan keempat, sebagaimana disorot pada Gambar~\ref{fig:pso-step45}, terdiri atas menteri pada (3,1) dan menteri pada (4,1) yang berada pada kolom 1. Sama halnya dengan pasangan pada Gambar~\ref{fig:pso-step43}, pasangan ini juga merupakan pelanggaran \textit{adjacency} (telah disorot pada Gambar~\ref{fig:pso-step39}) dan dihitung pada kedua komponen.

\begin{figure}[!htbp]
	\centering
	\includegraphics[width=0.27\linewidth]{./pso/45.png}
	\caption{Pasangan \textit{attacking} keempat: menteri pada (3,1) dan (4,1) berada pada kolom yang sama. Pasangan ini juga merupakan pelanggaran \textit{adjacency}.}
	\label{fig:pso-step45}
\end{figure}

Pasangan kelima, sebagaimana disorot pada Gambar~\ref{fig:pso-step46}, terdiri atas menteri pada (4,2) dan menteri pada (4,4) yang berada pada baris 4. Pasangan ini tidak bertetangga langsung sehingga hanya tergolong sebagai pelanggaran \textit{attacking}.

\begin{figure}[!htbp]
	\centering
	\includegraphics[width=0.27\linewidth]{./pso/43.png}
	\caption{Pasangan \textit{attacking} kelima: menteri pada (4,2) dan (4,4) berada pada baris yang sama.}
	\label{fig:pso-step46}
\end{figure}

Dengan kelima pasangan tersebut, nilai \textit{attackingViolation} adalah $5$. Nilai \textit{fitness} akhir partikel 2 setelah iterasi pertama dihitung berdasarkan Persamaan~\ref{eq:fitness-pso} dan menghasilkan nilai $11$, sebagaimana ditunjukkan pada Persamaan~\ref{eq:contoh-fitness-final}.

\begin{equation}
	fitness = w_1 \cdot \textit{attackingViolation} + w_2 \cdot \textit{adjacencyViolation} = (1 \cdot 5) + (2 \cdot 3) = 5 + 6 = 11
	\label{eq:contoh-fitness-final}
\end{equation}

Nilai \textit{fitness} ini lebih rendah dari \textit{fitness} awal partikel 2 (yang sebesar $20$), sehingga \textit{pBest} partikel 2 diperbarui dengan posisi baru. Selain itu, karena nilai ini ($11$) juga lebih rendah daripada \textit{fitness} \textit{nBest} sebelumnya (yaitu $14$ dari partikel 1), maka \textit{nBest} \textit{neighborhood} diperbarui menjadi partikel 2.

Proses pembaruan kecepatan, pembaruan posisi, dan evaluasi \textit{fitness} ini berulang pada iterasi-iterasi berikutnya. Apabila salah satu partikel mencapai \textit{fitness} bernilai 0 (yaitu konfigurasi tanpa pelanggaran), algoritma berhenti dan melaporkan partikel tersebut sebagai solusi valid. Apabila tidak, pencarian dilanjutkan hingga batas iterasi atau kondisi stagnasi tercapai.